% Electoral Equilibrium: Stochastic Coalition Optimization Under Political Shocks
% Juhi Damley, Gaston Espinosa
% Claremont McKenna College — SRP Summer 2026

\documentclass[12pt, letterpaper]{article}

% --- Packages ---
\usepackage[margin=1in]{geometry}
\usepackage{amsmath, amssymb, amsthm}
\usepackage{booktabs}
\usepackage{graphicx}
\usepackage{hyperref}
\usepackage{natbib}
\usepackage{setspace}
\usepackage{microtype}
\usepackage{xcolor}
\usepackage{algorithm}
\usepackage{algpseudocode}
\usepackage{array}
\usepackage{longtable}
\usepackage{caption}
\usepackage{subcaption}

\hypersetup{
  colorlinks = true,
  linkcolor  = blue!60!black,
  citecolor  = blue!60!black,
  urlcolor   = blue!60!black,
}

\doublespacing

% --- Theorem environments ---
\newtheorem{proposition}{Proposition}
\newtheorem{lemma}{Lemma}
\newtheorem{remark}{Remark}
\theoremstyle{definition}
\newtheorem{definition}{Definition}

% --- Math notation ---
\newcommand{\mueff}{\mu_{\mathrm{eff}}}
\newcommand{\Veq}{V_{\mathrm{eq}}}
\newcommand{\SigmaDelta}{\Sigma_{\Delta}}
\newcommand{\lambdaone}{\lambda_1}
\newcommand{\lambdatwo}{\lambda_2}
\newcommand{\lambdathree}{\lambda_3}
\newcommand{\Phicdf}{\Phi}
\newcommand{\boldw}{\mathbf{w}}
\newcommand{\boldmu}{\boldsymbol{\mu}}

% -----------------------------------------------------------------------
\begin{document}

\title{Electoral Equilibrium: A Stochastic Optimization Framework\\
       for Coalition Rebalancing Under Political Shocks}

\author{
  Juhi Damley \\
  Department of Mathematics \\
  Claremont McKenna College \\
  \texttt{jdamley@cmc.edu}
  \and
  Gaston Espinosa \\
  Department of Mathematics \\
  Claremont McKenna College \\
  \texttt{gespinosa@cmc.edu}
}

\date{June 2026}

\maketitle

\begin{abstract}
We present Electoral Equilibrium, a stochastic optimization framework that estimates
how a political party's voter coalition must rebalance after an exogenous shock.
Given a free-text description of a political event, the framework outputs
(i) updated coalition weights over five racial/ethnic blocs, and
(ii) a 90\% Monte Carlo interval on the post-shock probability of winning the
Electoral College. The framework combines three independently validated components:
a QLoRA-fine-tuned Mistral 7B language model (rank-16 adapter, MAE 0.036 on a
111-example held-out eval set) that translates shock narratives into per-stratum
loyalty-shift bins via constrained decoding, a disciplined quasi-convex program
(DQCP) that maximizes the probability of meeting an Electoral College threshold
under a Logistic-Normal ILR uncertainty model, and an NLP pipeline (RoBERTa +
SetFit) that constructs training data from social media and news archives.
The LLM is fine-tuned on a 1,136-record catalog of historical shock events
(577 hand-authored and 559 Gemini-reviewed synthetic examples validated by
maximum mean discrepancy, PCA alignment, and pairwise correlation difference
diagnostics).
The demographic architecture models three independent marginal strata---race,
religion, and gender---without nested cross-tabulations, acknowledging the
ecological fallacy as an explicit approximation. Iterative proportional fitting
(IPF) calibrates the stratum layer weights from 20 historical election cycles
(1948--2024); the raked solution assigns substantially higher explanatory weight
to gender ($\lambda_3 = 0.662$) than to race ($\lambda_1 = 0.114$), recovering
from first principles the documented growth of the gender gap since 1980.
A Gaussian Process leave-one-cycle-out classifier trained on panel vote shares
achieves 85\% accuracy and Brier score 0.071 on the 20-cycle panel.
Electoral College asymmetry analysis places the Democrat win threshold at
50.66\% of the two-party popular vote versus 49.34\% for Republicans;
this threshold is an empirical estimate subject to ongoing validation
(Section~\ref{sec:veq_results}).
\end{abstract}

\bigskip
\noindent\textbf{Keywords:} coalition optimization, stochastic programming, Gaussian
processes, electoral forecasting, demographic modeling, quasi-convex programming,
Logistic-Normal distribution

\newpage

% -----------------------------------------------------------------------
\section{Introduction}
\label{sec:intro}

Political parties do not operate in equilibrium. Exogenous shocks---Supreme Court
rulings, geopolitical crises, candidate scandals, policy reversals---shift the
loyalties of demographic subgroups, sometimes dramatically and rapidly. A party's
strategic challenge after such a shock is to determine how to rebalance its
coalition: which groups to mobilize, which to cede, and at what allocation of
organizing effort. This rebalancing problem is structurally similar to portfolio
optimization under uncertainty, where the ``assets'' are demographic blocs, the
``returns'' are vote shares, and the ``win condition'' is an Electoral College
threshold rather than a portfolio return target.

The analogy is not new. Erikson, MacKuen, and Stimson's \citeyear{EriksonEtAl2002}
\textit{The Macro Polity} established that presidential approval and economic
sentiment drive aggregate electoral outcomes, and a large subsequent literature has
modeled vote shares as functions of economic fundamentals
\citep{Abramowitz2008, Hibbs2012}. More recently, the growth of individual-level
survey data has enabled multilevel regression and poststratification (MRP) approaches
that decompose vote share across demographic subgroups \citep{GelmanLittle1997,
Park2004, Lax2009}. High-dimensional Bayesian hierarchical models now achieve
sub-percentage-point state-level forecast accuracy in U.S. presidential elections
\citep{Linzer2013, Economist2020}.

Two gaps in this literature motivate the present work. First, existing forecasting
frameworks are \textit{predictive}: they estimate outcomes given current conditions.
They do not answer the prescriptive question of how a party should rebalance its
coalition to improve electoral prospects after a shock shifts the baseline. Second,
the treatment of uncertainty in coalition-level models is typically frequentist---
Bayesian models exist but standard Dirichlet or normal approximations to the
coalition weight simplex distort the geometry of the uncertainty in ways that
matter for tail-event scenarios (wave elections, demographic realignments).

We address both gaps. The Electoral Equilibrium framework answers the prescriptive
question directly: given a shock and its estimated per-bloc loyalty shifts, which
reallocation of organizing effort maximizes the probability of meeting the Electoral
College win threshold? The answer is computed via a disciplined quasi-convex program
\citep{DiamondBoyd2019} that generalizes the classical Sharpe-ratio portfolio
objective to the simplex-constrained coalition setting. Uncertainty is propagated
via a Logistic-Normal distribution in isometric log-ratio (ILR) coordinates
\citep{Egozcue2003}, which correctly handles wave elections where all blocs move
together---a case that Dirichlet sampling cannot represent because its off-diagonal
covariances are by construction negative.

The framework operates as a three-stage pipeline:

\begin{enumerate}
  \item \textbf{Shock estimation (Stage~1):} A QLoRA-fine-tuned Mistral 7B language
        model \citep{Jiang2023, Hu2022} translates a free-text shock narrative
        into per-bloc, per-stratum loyalty-shift bins using constrained decoding
        \citep{Willard2023}.

  \item \textbf{Coalition optimization (Stage~2):} A CVXPY quasi-convex program
        maximizes $\Phi\!\left(\bigl(\mueff(\boldw) - \Veq\bigr) /
        \sqrt{\lambdaone^2\, \boldw^\top \SigmaDelta \boldw}\right)$
        over the probability simplex, where $\Phi$ is the standard normal CDF,
        $\mueff$ is the three-stratum effective loyalty scalar, $\Veq$ is the
        EC-adjusted win threshold, and $\SigmaDelta$ is the empirical covariance
        of race-bloc loyalty shifts.

  \item \textbf{NLP pipeline (Stage~3):} A RoBERTa sentence encoder \citep{Liu2019}
        combined with a SetFit few-shot demographic classifier \citep{Tunstall2022}
        produces bloc-weighted sentiment scores from social media and news archives,
        supplying fine-tuning training data for Stage~1.
\end{enumerate}

Our primary contributions are:

\begin{enumerate}
  \item A three-stratum additive independence architecture (race, religion, gender)
        calibrated by iterative proportional fitting, with explicit ecological
        fallacy acknowledgment and an empirical comparison of additive versus
        raked outputs.

  \item An Electoral College--adjusted win threshold derived from logistic regression
        on 20 historical cycles, placing the Democrat threshold at $\Veq = 0.5066$
        versus $\Veq = 0.4934$ for Republicans.

  \item A quasi-convexity proof for the Sharpe-ratio coalition objective and a
        demonstration that min-variance optimization---the standard portfolio
        baseline---systematically assigns zero weight to Latino and Asian blocs
        while the DQCP optimizer restores positive weights in deficit scenarios.

  \item Empirical validation of the GP baseline classifier (85\% LOCO accuracy,
        Brier 0.071) and documentation of three data-quality issues---Perot
        three-party contamination, ANES oversampling in 1988, and pre-realignment
        era drift---whose corrections non-trivially affect downstream estimates.

  \item Fine-tuning results for the Stage~1 Mistral 7B QLoRA adapter: rank-16
        achieves MAE 0.036 on a held-out eval set of 111 synthetic examples,
        outperforming rank-32 (MAE 0.037, 2.8\% worse) despite being trained on
        a smaller parameter budget. We document a directional error on the Newsom
        2028 nomination scenario---the pre-augmentation model incorrectly predicts
        \texttt{slight\_pos} for all minority blocs---and report the retraining
        procedure on a 1,136-record augmented dataset designed to correct it.
\end{enumerate}

The remainder of the paper proceeds as follows. Section~\ref{sec:background}
reviews related work. Section~\ref{sec:framework} describes the framework
architecture. Section~\ref{sec:data} details the panel construction. Section~\ref{sec:methods}
presents the methodology for each component. Section~\ref{sec:results} reports
empirical findings. Section~\ref{sec:discussion} discusses limitations, the
ecological fallacy, and calibration against prediction markets. Section~\ref{sec:conclusion}
concludes.

% -----------------------------------------------------------------------
\section{Background and Related Work}
\label{sec:background}

\subsection{Electoral forecasting}

Quantitative electoral forecasting in the United States divides broadly into
structural models, which predict outcomes from macroeconomic and political
``fundamentals'' \citep{Fair1978, Abramowitz2008, Hibbs2012}, and
polling-aggregation models, which update structural priors with survey data
\citep{Erikson2014, Linzer2013}. A third class uses individual-level panel
surveys via MRP to produce small-area estimates \citep{GelmanLittle1997,
Lax2009, Wang2015}. All three classes are primarily predictive; none frames
the coalition rebalancing question as an optimization problem.

The closest antecedents to our work are the ``margin of error'' decomposition
models \citep{King1991} and simulation-based approaches to electoral uncertainty
\citep{Jackman2005, Strauss2007}. Gaussian Process models for vote shares appear
in \citet{Linzer2013} and \citet{Lock2010}, who use GP priors over election-day
outcomes. Our use of GPR as a baseline win-probability classifier over 20-cycle
panel data differs in purpose: we are not forecasting an upcoming election but
validating that our panel vote-share features are informative about realized
outcomes, which is the necessary precondition for the optimizer's objective.

\subsection{Portfolio theory and coalition modeling}

The analogy between political coalitions and investment portfolios has been noted
informally by practitioners but formally developed by only a handful of political
economists. \citet{Enelow1984} established spatial models of coalition building.
Closer to our framing, \citet{Erikson2012} used a variance-minimization framing
to model party positioning. The min-variance interpretation fails in exactly the
deficit scenarios we target: a party whose effective loyalty is below the win
threshold cannot improve its probability of winning by minimizing variance---it
must maximize the probability of the tail event \citep{Roy1952}.

The Sharpe-ratio objective we employ has antecedents in financial portfolio theory
\citep{Sharpe1966} and in reliability engineering \citep{Roy1952}. Its adaptation
to the quasi-convex programming setting via DQCP is due to \citet{DiamondBoyd2019},
who prove that a ratio of an affine function to a positive convex function is
quasi-convex and can be solved exactly by disciplined convex programming software
such as CVXPY.

\subsection{Compositional data and the simplex}

Coalition weight vectors lie on the probability simplex. Standard normal or
Dirichlet approximations distort the geometry of uncertainty in compositional
data in well-documented ways \citep{Aitchison1986}. The Dirichlet distribution
forces all off-diagonal covariances negative, making it impossible to model the
positive correlations that characterize wave elections. The delta method with
a floor constraint distorts the Aitchison geometry non-linearly near the boundary
\citep{Egozcue2003}.

The correct approach is the isometric log-ratio (ILR) transformation with
Helmert contrast matrices, which maps the interior of the simplex to an
unconstrained $\mathbb{R}^{K-1}$ space where Gaussian distributions are
appropriate \citep{Pawlowsky2015}. We use Logistic-Normal ILR draws to
propagate the race-bloc covariance $\SigmaDelta$ through the Monte Carlo,
producing the 90\% credible interval on $P(\text{win})$.

\subsection{Demographic modeling and ecological inference}

The ecological fallacy---inferring individual-level behavior from group-level
aggregates---is central to demographic electoral modeling
\citep{Robinson1950, King1997}. MRP addresses this by modeling individual vote
probabilities as a function of individual-level covariates, then integrating
over the population distribution \citep{Park2004}. We do not use MRP: it
requires full joint cross-tabulations (race $\times$ religion $\times$ gender)
that produce sparse cells at the intersection of minority groups, and its
implementation is outside the scope of this project's timeline. Instead, we
use three parallel marginal strata with an additive independence approximation
and compare its outputs to an IPF-calibrated raked alternative.

% -----------------------------------------------------------------------
\section{Framework Architecture}
\label{sec:framework}

\subsection{Overview}

The Electoral Equilibrium framework is a three-stage stochastic pipeline.
Stage~1 translates a user-supplied shock narrative into per-stratum loyalty
shift estimates. Stage~2 finds the coalition allocation that maximizes win
probability under those shifts. Stage~3 continuously generates training data
for Stage~1 from social media and news sources.

Each stage produces a frozen artifact (\texttt{ShockResponseData},
\texttt{EquilibriumData}, \texttt{SimulationData}) written to disk in a
deterministic format keyed by a global random seed. All stochastic operations
derive their seeds from a single top-level seed via $\texttt{derive\_seed}(s_0,
\text{stage\_name})$, ensuring byte-identical reproducibility across runs.

\subsection{Demographic architecture: three parallel strata}

The U.S. electorate is modeled as three independent marginal strata, each
covering approximately 100\% of voters:

\begin{itemize}
  \item \textbf{Stratum~1 (Race/Ethnicity):} Five blocs---African American
    ($\approx$12\%), Latino ($\approx$11\%), Asian ($\approx$5\%), White
    ($\approx$62\%), and Other Race ($\approx$10\%)---constitute the
    optimizer's primary decision variables $\boldw \in \Delta^4$ (the
    4-simplex).

  \item \textbf{Stratum~2 (Religion):} Seven blocs---Evangelical
    ($\approx$24\%), Catholic ($\approx$21\%), Protestant ($\approx$13\%),
    Secular ($\approx$26\%), Jewish ($\approx$2\%), Muslim ($\approx$1\%),
    and Other Religion ($\approx$13\%)---with fixed population weights
    $\mathbf{v}$ drawn from the voter panel.

  \item \textbf{Stratum~3 (Gender):} Three blocs---Women ($\approx$52\%),
    Men ($\approx$47\%), and Other Gender ($\approx$1\%)---with fixed
    population weights $\mathbf{g}$.
\end{itemize}

The three strata are explicitly not nested: we do not require race $\times$
religion $\times$ gender cross-tabulations. This avoids sparse-cell problems
(e.g., Asian $\times$ Muslim $\times$ Other Gender $\approx 0.05\%$ of the
electorate) at the cost of an additive independence approximation discussed
in Section~\ref{sec:eco_fallacy}.

\subsection{Effective loyalty scalar}
\label{sec:mueff}

The three-stratum effective loyalty scalar is:

\begin{equation}
  \mueff(\boldw) \;=\;
    \lambdaone \sum_{i} w_i\, \mu_i^{(\text{race})}
  \;+\; \lambdatwo \sum_{r} v_r\, \mu_r^{(\text{rel})}
  \;+\; \lambdathree \sum_{g} g_g\, \mu_g^{(\text{gen})},
  \label{eq:mueff}
\end{equation}

where $\mu_i^{(\text{race})}$, $\mu_r^{(\text{rel})}$, $\mu_g^{(\text{gen})}$
are the within-bloc Democratic vote shares estimated from the voter panel, and
$\lambdaone + \lambdatwo + \lambdathree = 1$ with $\lambda_k \geq 0$. The
optimizer controls only $\boldw$; the religion and gender weights $\mathbf{v}$
and $\mathbf{g}$ are fixed from the panel and are not decision variables.

The layer weights $\boldsymbol{\lambda} = (\lambdaone, \lambdatwo, \lambdathree)$
are calibrated from historical data by iterative proportional fitting
(Section~\ref{sec:raking}). We maintain two calibrations in parallel: an
\textit{additive prior} (50/30/20) reflecting conventional wisdom about race's
primacy in U.S. electoral politics, and a \textit{raked} calibration
(11.4/22.4/66.2) derived from historical two-party vote shares. The divergence
between the two is itself a substantive finding (Section~\ref{sec:raking_results}).

\subsection{Win condition and Electoral College adjustment}
\label{sec:veq}

The optimizer targets $\mueff(\boldw) \geq \Veq$, where $\Veq$ is the
popular-vote two-party share at which $P(\text{party wins EC}) = 0.50$. We
derive $\Veq$ by logistic regression:

\begin{equation}
  P(\text{Dem EC win}) \;=\;
    \sigma\!\left(A \cdot x_{\text{dem2p}} + B\right),
  \qquad
  \Veq^{(\text{Dem})} = -B/A,
  \label{eq:veq}
\end{equation}

where $x_{\text{dem2p}}$ is the Democratic two-party popular vote share and
$\sigma$ is the logistic function. This is estimated on all 20 historical
cycles, 1948--2024; see Section~\ref{sec:veq_results}.

% -----------------------------------------------------------------------
\section{Data}
\label{sec:data}

\subsection{Survey sources}

The voter panel integrates six data sources across three strata:

\begin{enumerate}
  \item \textbf{ANES Cumulative Data File (1948--2020):} American National
    Election Studies provides the core longitudinal panel for race, gender,
    and some religion variables. We use the VCF-coded cumulative file with
    auto-extracted label dictionary from the codebook PDF.

  \item \textbf{Cooperative Election Study (CES, 2006--2024):} Over 700,000
    respondents across six election cycles. CES is the primary source for
    religion blocs post-2006 and provides evangelical/Protestant
    disaggregation via the \texttt{race\_h} Hispanic-override variable and
    evangelical flag column.

  \item \textbf{General Social Survey (GSS, 1972--2022):} Primary source for
    religion bloc time series. Weights use the \texttt{wtssnrps} $\to$
    \texttt{wtssps} $\to$ \texttt{wtss} fallback chain.

  \item \textbf{National Election Pool Exit Polls (NEP, 2000--2024):} Gold
    standard for election-night demographic vote shares. Used as the primary
    source for race blocs in cycles where available and for benchmarking all
    panel estimates.

  \item \textbf{NORC Public Opinion Research Survey (NPORS, 2024):} Race,
    religion, and gender vote shares for the 2024 cycle.

  \item \textbf{Democracy Fund VOTER Panel:} Longitudinal voter tracking
    across multiple cycles, providing panel continuity.
\end{enumerate}

\subsection{Panel construction and conflict resolution}

The panel covers 20 presidential election cycles (1948--2024) across 15
canonical demographic blocs (five race, seven religion, three gender), for
a maximum of 300 bloc-cycle cells. After data collection, coverage was 234
cells (78.0\%); all 15 blocs are present in every cycle from 2004 to 2024
(100\% modern coverage).

When multiple sources report a vote share for the same bloc-cycle cell,
conflicts are resolved by inverse-standard-error weighting:
\begin{equation}
  \hat{v}_{\text{merged}} = \frac{\sum_s w_s\, v_s}{\sum_s w_s},
  \qquad
  w_s = \frac{1}{\text{SE}_s},
  \label{eq:merge}
\end{equation}
with SE estimated from the Kish effective sample size for ANES/GSS/CES and
from the binomial formula for NEP. Each conflict is logged with per-source
values.

\subsection{Three-party contamination correction}
\label{sec:perot}

Survey panels record the raw Democratic fraction of all votes cast, including
third-party votes. In 1992, Ross Perot received 18.9\% of the popular vote;
Clinton's raw Democratic fraction ($\approx$43\%) substantially understates
his two-party share (53.5\%). We correct each 1992 bloc using bloc-specific
Perot support from the 1992 National Exit Poll:

\begin{equation}
  v_{i,1992}^{(2p)} = \frac{v_{i,1992}^{(\text{raw})}}{1 - p_{i,1992}},
  \label{eq:perot}
\end{equation}

where $p_{i,1992}$ is the Perot share for bloc $i$ (White: 21\%, African
American: 7\%, Latino: 14\%, etc.). Without this correction, the 1992
panel row is misclassified as a Republican win and the covariance matrix
$\SigmaDelta$ carries inflated race-bloc variance.

\subsection{Structural gaps and imputation}

Evangelical and secular religion blocs lack reliable data before 2000, when
the relevant institutions (Moral Majority, ``nones'') either did not exist or
were insufficiently represented in survey instruments. Muslim and Asian blocs
are similarly sparse before the 1965 Immigration Act. We do not impute these
structural gaps; instead, LLM fine-tuning for these blocs is restricted to
2004--2024 cycles where data exists. Three isolated cells (other\_gender for
five cycles, Muslim 2004, other\_race 1948) are imputed from documented
external sources with explicit source tags.

\subsection{Panel quality benchmarks}

Table~\ref{tab:benchmarks} reports estimated Democratic vote shares from the
panel alongside NEP and Pew benchmarks for the 13 Democratic winning cycles.
All race blocs are within 7 percentage points of benchmarks; notable
discrepancies are historically explained rather than indicative of data
errors (see Section~\ref{sec:discussion}).

\begin{table}[h]
\centering
\caption{Panel vote-share estimates vs.\ benchmark sources (Democratic winning cycles, 1948--2024).}
\label{tab:benchmarks}
\begin{tabular}{lcccl}
\toprule
Bloc & Panel $\mu$ & Benchmark & Deviation & Notes \\
\midrule
African American & 0.883 & 0.870 & $+$1.3 pp & NEP \\
Latino           & 0.663 & 0.620 & $+$4.3 pp & NEP \\
Asian            & 0.610 & 0.660 & $-$5.0 pp & AAPI Data \\
White            & 0.456 & 0.390 & $+$6.6 pp & NEP; ANES 2020 oversample \\
Other Race       & 0.644 & 0.580 & $+$6.4 pp & NEP; heterogeneous bloc \\
\midrule
Evangelical      & 0.275 & 0.230 & $+$4.5 pp & Pew; pre-Moral-Majority cycles \\
Catholic         & 0.568 & 0.520 & $+$4.8 pp & NEP \\
Secular          & 0.703 & 0.680 & $+$2.3 pp & Pew \\
Jewish           & 0.789 & 0.720 & $+$6.9 pp & Pew; CES education skew \\
\midrule
Women            & 0.534 & 0.540 & $-$0.6 pp & NEP \\
Men              & 0.487 & 0.430 & $+$5.7 pp & NEP; pre-1980 gender gap \\
\bottomrule
\end{tabular}
\end{table}

% -----------------------------------------------------------------------
\section{Methodology}
\label{sec:methods}

\subsection{Gaussian Process baseline classifier}
\label{sec:gp}

The GP classifier validates that panel vote-share features are informative
about realized election outcomes and provides a calibrated prior on
$P(\text{win} \mid \text{coalition})$ for the optimizer.

\subsubsection{Model specification}

We fit a \texttt{GaussianProcessRegressor} on binary $\{0,1\}$ outcomes using
an additive kernel $k = k_{\mathrm{RBF}} + k_{\mathrm{Mat\'ern}(\nu=1.5)}$.
We use GPR rather than GPC because the paper requires cycle-level posterior
standard deviations $\sigma_c = \mathrm{std}[f(x_c)]$ as epistemic uncertainty
estimates. The sklearn GPC (Laplace approximation) does not expose the latent
posterior standard deviation through its public API; GPR returns it directly
via \texttt{predict(return\_std=True)}. GPR on binary labels is technically
misspecified but manages this via:

\begin{itemize}
  \item Predictions clipped to $[0,1]$.
  \item Observation noise $\alpha = 0.10$, which models $\approx$3 pp survey
        measurement error in standardized label space and prevents exact
        interpolation through training points.
  \item Feature standardization via \texttt{StandardScaler} fitted on all
        available cycles before the leave-one-cycle-out (LOCO) loop.
\end{itemize}

Without \texttt{StandardScaler}, raw vote-share features (span $\approx$0.3--0.9)
produce kernel length-scale collapse: the optimizer drives both RBF and
Matérn length scales to their lower bounds ($10^{-5}$), reducing the model
to a near-delta function and producing coin-flip accuracy.

\subsubsection{Feature construction}

Each training cycle is represented by a $d = 9$-dimensional feature vector:
\begin{itemize}
  \item Five race-bloc Democratic vote shares $v_{i,c}^{(\text{race})}$.
  \item Three stratum-weighted coalition averages:
    \begin{align}
      \bar{\mu}_c^{(\text{race})} &=
        \textstyle\sum_i \tilde{e}_i^{(\text{race})} \, v_{i,c}^{(\text{race})},
      \label{eq:murace} \\
      \bar{\mu}_c^{(\text{rel})} &=
        \textstyle\sum_r \tilde{e}_r^{(\text{rel})} \, v_{r,c}^{(\text{rel})},
      \notag \\
      \bar{\mu}_c^{(\text{gen})} &=
        \textstyle\sum_g \tilde{e}_g^{(\text{gen})} \, v_{g,c}^{(\text{gen})},
      \notag
    \end{align}
    where $\tilde{e}$ are approximate national electorate shares normalized
    over blocs present in cycle $c$.
  \item A normalized cycle-year feature $t_c = (c - 1948) / 76 \in [0, 1]$,
        which gives the GP temporal ordering so that pre-realignment cycles
        (1948--1968) receive lower covariance with modern cycles.
\end{itemize}

The stratum-weighted averages were added after finding that the GP, presented
with only five race-bloc features, fixated on African American loyalty in 2024
($\mu_{\text{AA},2024} = 0.838$, matching a 2020 Democratic win pattern) while
missing the coalition-wide deficit: $\bar{\mu}_{2024}^{(\text{race})} = 0.486$,
$\bar{\mu}_{2024}^{(\text{rel})} = 0.439$, $\bar{\mu}_{2024}^{(\text{gen})} = 0.483$,
all below $\Veq = 0.507$.

\subsubsection{Leave-one-cycle-out validation}
\label{sec:loco}

LOCO-CV trains on 19 of 20 cycles and predicts the held-out cycle. For each
fold $c$, we refit the scaler and GP on the training set and report
\texttt{prob\_win}$_c$ and \texttt{prob\_std}$_c$ (the GP posterior standard
deviation) on the held-out cycle.

\subsection{Moment estimation}
\label{sec:moments}

The optimizer requires $\boldmu^{(P)}$ (mean within-bloc Democratic vote share
across winning cycles for party $P$) and $\SigmaDelta$ (the $5\times5$
empirical race-bloc covariance across all cycles).

\textbf{Winning cycles} are identified from the external ground-truth table
\texttt{\_PRES\_DEM\_2P\_SHARE}, not from the panel-derived vote-share
heuristic that fails in three-party years. The panel was verified to agree
with certified two-party returns on 16 of 20 cycles; the four exceptions
(1952, 1988, 1992, 2004) are corrected by the external table.

\textbf{Covariance} $\SigmaDelta$ is estimated from all 20 cycles (not
conditioned on winning) using the sample covariance with $\text{ddof} = 1$.
Ledoit-Wolf shrinkage \citep{Ledoit2004} is applied to guarantee
positive-definiteness when the training set for a particular LOCO fold
has fewer than $5 + 1$ observations:

\begin{equation}
  \hat{\Sigma}_{\mathrm{LW}} = (1 - \alpha)\, \hat{\Sigma}_{\mathrm{sample}}
  + \alpha\, \hat{\mu}_{\mathrm{sample}} I_{5\times5}.
  \label{eq:ledoitwolf}
\end{equation}

\subsection{Layer weight calibration by iterative proportional fitting}
\label{sec:raking}

We calibrate $\boldsymbol{\lambda} = (\lambdaone, \lambdatwo, \lambdathree)$
by minimizing the MSE of $\mueff$ against historical two-party Democratic vote
shares:

\begin{equation}
  \hat{\boldsymbol{\lambda}} = \operatorname*{arg\,min}_{\boldsymbol{\lambda} \in \Delta^2}
  \sum_{c} \left(
    \lambdaone\, \bar{\mu}_c^{(\text{race})}
    + \lambdatwo\, \bar{\mu}_c^{(\text{rel})}
    + \lambdathree\, \bar{\mu}_c^{(\text{gen})}
    - y_c
  \right)^2,
  \label{eq:raking_obj}
\end{equation}

where $y_c$ is the historical two-party Democratic popular vote share and
$\Delta^2 = \{\boldsymbol{\lambda} : \sum_k \lambda_k = 1,\ \lambda_k \geq 0\}$.

The three stratum series $\bar{\mu}^{(\text{race})}, \bar{\mu}^{(\text{rel})},
\bar{\mu}^{(\text{gen})}$ are nearly collinear---all track the partisan tide
with pairwise correlations $\approx 0.92$---making the unregularized objective
nearly flat and causing first-order methods to converge extremely slowly:
multiplicative IPF required $\approx$288 iterations; projected gradient descent
required $\approx$101 iterations; FISTA (Nesterov momentum) required
$\approx$35 iterations.

The correct approach is Tikhonov regularization toward an equal-weight prior:

\begin{equation}
  \left(\mathbf{S}^\top\mathbf{S} + \rho I\right)\, \boldsymbol{\lambda}
  = \mathbf{S}^\top \mathbf{y} + \rho\, \boldsymbol{\lambda}_{\text{prior}},
  \label{eq:tikhonov}
\end{equation}

where $\mathbf{S} \in \mathbb{R}^{n \times 3}$ stacks the stratum series,
$\rho = 0.01$ is the regularization weight, and
$\boldsymbol{\lambda}_{\text{prior}} = (1/3, 1/3, 1/3)$ is the equal-weight
prior. The 1\% regularization reduces the condition number of
$\mathbf{S}^\top\mathbf{S}$ from $\approx$120 to $\approx$12, allowing an
exact one-shot solve followed by the Duchi et al.\ \citeyear{Duchi2008} simplex
projection. An active-set outer loop re-solves on the reduced simplex face
if the projection zeros out a stratum; this typically requires one additional
pass. The total iteration count is two.

\subsection{Coalition optimizer (DQCP)}
\label{sec:optimizer}

\subsubsection{Objective formulation}

The post-shock effective loyalty scalar is:

\begin{equation}
  \tilde{\mu}_{\mathrm{eff}}(\boldw) \;=\;
    \lambdaone \sum_i w_i\,(\mu_i + \delta_i^{(\text{race})})
  \;+\; \lambdatwo \sum_r v_r\,(\mu_r + \delta_r^{(\text{rel})})
  \;+\; \lambdathree \sum_g g_g\,(\mu_g + \delta_g^{(\text{gen})}),
  \label{eq:mueff_shock}
\end{equation}

where $\boldsymbol{\delta}$ are the shock deltas from Stage~1.

The optimizer maximizes the probability that $\tilde{\mu}_{\mathrm{eff}}(\boldw)$
exceeds $\Veq$ under a Gaussian model for the race-level shock uncertainty:

\begin{equation}
  \max_{\boldw \in \Delta^4}
  \Phi\!\left(
    \frac{\tilde{\mu}_{\mathrm{eff}}(\boldw) - \Veq}
         {\sqrt{\lambdaone^2\, \boldw^\top \SigmaDelta \boldw}}
  \right).
  \label{eq:dqcp_obj}
\end{equation}

Since $\Phi$ is monotone increasing, this is equivalent to maximizing the
argument. We prove below that the argument is quasi-convex in $\boldw$.

\subsubsection{Quasi-convexity proof}

\begin{proposition}
\label{prop:qcvx}
The objective in \eqref{eq:dqcp_obj} is quasi-convex in $\boldw \in \Delta^4$.
\end{proposition}

\begin{proof}
Let $f(\boldw) = g(\boldw) / h(\boldw)$ where
$g(\boldw) = \tilde{\mu}_{\mathrm{eff}}(\boldw) - \Veq$ and
$h(\boldw) = \sqrt{\lambdaone^2\, \boldw^\top \SigmaDelta \boldw}$.

\textit{Numerator.} $g(\boldw)$ is affine in $\boldw$ by inspection of
\eqref{eq:mueff_shock}: it equals $\lambdaone\, \boldmu^{(\text{race})\top}\boldw$
plus terms that do not depend on $\boldw$. An affine function is simultaneously
convex and concave.

\textit{Denominator.} $\SigmaDelta$ is positive semi-definite by construction;
Ledoit-Wolf shrinkage guarantees strict positive-definiteness. Therefore
$\boldw^\top \SigmaDelta \boldw$ is a strictly convex quadratic form for
$\boldw \in \text{int}(\Delta^4)$, and $h(\boldw) = \lambdaone
\sqrt{\boldw^\top \SigmaDelta \boldw}$ is convex (composition of a convex
function with $\sqrt{\cdot}$, which is concave and monotone increasing on
$\mathbb{R}_+$, preserves convexity). Moreover $h(\boldw) > 0$ for all
$\boldw$ with $w_i > 0$ for some $i$.

\textit{Ratio.} By Proposition~1 of \citet{DiamondBoyd2019}, a ratio of an
affine function to a positive convex function is quasi-convex. Therefore
$f(\boldw)$ is quasi-convex on $\Delta^4$.

\textit{Composition.} $\Phi$ is monotone increasing, and the composition of a
quasi-convex function with a monotone increasing function is quasi-convex.
Therefore $\Phi(f(\boldw))$ is quasi-convex, and its maximization is a
quasi-convex program (QCP). \qed
\end{proof}

The program is solved by CVXPY with \texttt{problem.solve(qcp=True)}; the
implementation asserts \texttt{problem.is\_dcp(qcp=True)} at runtime as a
correctness guard.

\textbf{Why not min-variance?} The standard mean-variance portfolio objective
$\min_\boldw \boldw^\top \SigmaDelta \boldw$ subject to $\boldmu^\top \boldw
\geq \Veq$ is equivalent to asking: ``minimize uncertainty given the win
constraint.'' When $\boldmu^\top \boldw$ already exceeds $\Veq$ for the
high-loyalty blocs, this approach concentrates weight there and assigns zero
weight to blocs that add variance without a proportional loyalty gain.
In our Democrat baseline, this produces $w_{\text{Latino}} = w_{\text{Asian}} = 0$
even though both groups represent substantial electorates. The quasi-convex
objective---maximizing $P(\text{win})$ rather than minimizing variance---restores
positive weight to these blocs in deficit scenarios where the path to $\Veq$
requires accepting additional variance.

\subsection{Monte Carlo uncertainty quantification}
\label{sec:mc}

The 90\% credible interval on $P(\text{win})$ is computed via 10,000
Logistic-Normal ILR draws. Let $\boldw^*$ be the optimizer solution.

\begin{enumerate}
  \item Map $\boldw^*$ to ILR coordinates: $\mathbf{z}^* = \mathbf{V}^\top
        \log(\boldw^*)$, where $\mathbf{V} \in \mathbb{R}^{5\times4}$ is the
        Helmert contrast matrix.
  \item Propagate $\SigmaDelta$ to ILR space via the Jacobian:
        $\Sigma_{\mathrm{ILR}} = J\, \SigmaDelta\, J^\top$.
  \item Draw $\mathbf{y}^{(n)} \sim \mathcal{N}(\mathbf{z}^*, \Sigma_{\mathrm{ILR}})$
        for $n = 1, \ldots, 10{,}000$.
  \item Back-transform: $\boldw^{(n)} = \mathrm{softmax}(\mathbf{V}\mathbf{y}^{(n)})$.
  \item Compute the win flag $\mathbf{1}[\mueff(\boldw^{(n)}) \geq \Veq]$ and
        the 5th and 95th percentiles as the 90\% CI bounds.
\end{enumerate}

The Dirichlet distribution was explicitly rejected for this step because its
covariance structure forces $\mathrm{Cov}(w_i, w_j) < 0$ for all $i \neq j$,
making it impossible to model wave elections where all blocs shift together.
The delta method with a floor at $w_i \geq 0.01$ was rejected because the
perturbation is not uniform in Aitchison geometry.

\begin{remark}[Scope of the reported CI]
The 90\% interval reported here reflects \emph{Monte Carlo sampling variance}:
how precisely we estimate the win fraction across $N = 10{,}000$ draws of
$\boldw^{(n)}$. This quantity shrinks as $1/\sqrt{N}$ and is not scientifically
meaningful as a measure of model uncertainty---it is a parameter we control, not
a property of the data. A scientifically meaningful interval would additionally
propagate uncertainty in the covariance estimate $\SigmaDelta$ itself, for
example by drawing $\hat{\Sigma} \sim \mathrm{Wishart}(\hat{\Sigma}_{\mathrm{LW}},
n_{\mathrm{cycles}})$ and repeating the full ILR Monte Carlo for each draw.
This outer loop is deferred to future work. In results tables and figures, the
reported CI should be read as ``simulation precision'' rather than
``epistemic uncertainty about $P(\text{win})$.''
\end{remark}

\subsection{Stage 1: LLM fine-tuning for shock estimation}
\label{sec:llm}

Stage~1 translates a user-supplied political shock narrative into per-stratum
loyalty shift bins. We fine-tune Mistral 7B Instruct v0.3 \citep{Jiang2023}
via QLoRA (rank~16, $\alpha = 32$, 4-bit NF4 quantization) on a catalog of
historical shock events.

\textbf{Output vocabulary.} The model is constrained to produce exactly one of
nine tokens per bloc:
$\{$\texttt{strong\_neg}, \texttt{mod\_neg}, \texttt{mild\_neg},
\texttt{slight\_neg}, \texttt{neutral}, \texttt{slight\_pos},
\texttt{mild\_pos}, \texttt{mod\_pos}, \texttt{strong\_pos}$\}$,
implemented via the \texttt{outlines} constrained decoding library
\citep{Willard2023}. Token midpoints span $[-0.12, +0.12]$ in vote-share
space. Per-bin uncertainty $\sigma_b$ is empirically estimated from
historical variation and stored in \texttt{configs/bin\_uncertainty.json};
conservative priors are used for bins with fewer than five historical examples.

\textbf{Training data.} Historical training records pair shock descriptions with
per-stratum delta bins derived from pre/post-event changes in survey vote
shares. The current training set contains \textbf{1,136 records}: 577 hand-authored
from historical shock events and 559 Gemini-reviewed synthetic examples. Synthetic
records are generated via a three-stage pipeline: (i) DeepSeek constrained to
the four-axis taxonomy (stratum $\times$ valence $\times$ intensity $\times$ party),
(ii) Gemini 2.5 Flash plausibility review (APPROVE / REVISE / HUMAN\_REVIEW
verdicts), and (iii) a mandatory 5\% human spot-check. Every synthetic row carries
a \texttt{source} provenance tag and \texttt{\_seed\_meta} block linking it to its
generating event and party. An additional 647 records flagged for human review are
\emph{excluded} from the current training set; their inclusion is contingent on
directional validation of the retrained model (see below).

Synthetic batches are validated before inclusion by three diagnostics:
(i) maximum mean discrepancy (MMD) with RBF kernel and the median heuristic
$\lambda = 1/(2\sigma^2)$ \citep{Gretton2012}; (ii) PCA alignment requiring
$>70\%$ of synthetic variance explained by the top-2 real principal components;
(iii) pairwise correlation difference (PCD) Frobenius norm. Per-row soft weights
$w = 0.5 \cdot \exp(-\lambda \cdot \mathrm{MMD}^2)$ downweight synthetic records
with poor distributional alignment; hard rejection is not applied.

\textbf{Fine-tuning results.} Table~\ref{tab:lora_ranks} reports evaluation
metrics for rank-16 and rank-32 adapters trained on an 80/20 stratified
train/eval split of the pre-augmentation 557-example dataset (446 train /
111 eval), evaluated by mean absolute error in delta-bin units and direction
accuracy (sign match).

\begin{table}[h]
\centering
\caption{QLoRA adapter rank comparison (Mistral 7B Instruct v0.3, 111-example
  held-out eval set stratified by party).}
\label{tab:lora_ranks}
\begin{tabular}{lcccc}
\toprule
Rank & Overall MAE & Direction acc.\ & Training loss & Adapter size \\
\midrule
r=16 & \textbf{0.0362} & 110/111 & 0.2331 & $\sim$84M params \\
r=32 & 0.0372          & 111/111 & 0.2693 & $\sim$168M params \\
\bottomrule
\end{tabular}
\end{table}

Rank-32 is 2.8\% worse in MAE and shows higher training loss, consistent with
overfitting on the 446-example training set. The rank-16 adapter at
\texttt{models/mistral-r16-hopper} is selected as the production adapter.
Highest-error blocs are \texttt{other\_gender} (MAE 0.049), \texttt{evangelical}
(0.045), and \texttt{african\_american} (0.042)---all blocs with sparse
synthetic training coverage that will improve when real RoBERTa-scored data
enters the training mix.

\textbf{Directional validation.} Running the pre-augmentation adapter on the
Newsom 2028 presidential nomination scenario produced \texttt{slight\_pos} for
all five race blocs---a politically incorrect prediction, as Newsom is expected
to hurt with White working-class and some Latino voters while possibly helping
with African American and secular blocs. The directional error indicates that
the model is biased toward positive bins by the synthetic training distribution
rather than learning genuinely event-specific partisan dynamics. The 1,136-record
augmented dataset is intended to correct this by expanding coverage of events
with mixed or negative minority-bloc effects. Validation of the retrained model
on the Newsom scenario is a required gate before the paper-baseline run.

\subsection{Stage 3: NLP pipeline}
\label{sec:nlp}

Stage~3 generates training data for Stage~1 from social media and news archives.
Author demographic inference uses a three-stage pipeline: (i) a SetFit
few-shot classifier \citep{Tunstall2022} on bio text, served from a
Raspberry Pi 5 running \texttt{all-MiniLM-L6-v2} via sentence-transformers
in CPU mode ($\approx$20\,ms per bio; Hailo NPU compilation deferred
pending vendor SDK availability); (ii) a keyword lexicon fallback
(\texttt{configs/race\_lexicon.json}, \texttt{religion\_lexicon.json},
\texttt{gender\_lexicon.json}); and (iii) a language-prior fallback based
on Pew Research population estimates. Posts assigned via the language prior
are excluded from both the mean $\boldmu$ and covariance $\SigmaDelta$
estimation---including them in only one would bifurcate the data-generating
process and violate the mean-variance optimizer's assumptions.

Sentiment toward each demographic bloc is scored by a bloc-weighted RoBERTa
sentence encoder \citep{Liu2019}. The resulting ElasticityScore scalars
$\in [-1, 1]$ serve as input features to the Stage~1 LLM alongside the
free-text shock description.

\subsubsection{Bio classifier validation}
\label{sec:bio_val}

The SetFit classifier was trained on 2,191 hand-labeled bios drawn from the
Bluesky and Reddit archives, split 80/20 train/eval.  Each bio received a
single majority label per stratum, assigned by keyword heuristics followed by
manual review.  Table~\ref{tab:setfit_f1} reports macro-averaged F1 on the
held-out eval split.

\begin{table}[h]
\centering
\caption{SetFit bio classifier validation results (macro-F1, 438-bio eval split).}
\label{tab:setfit_f1}
\begin{tabular}{lcc}
\toprule
Stratum & Classes & Macro-F1 \\
\midrule
Race/ethnicity  & 5 & 0.97 \\
Religion        & 7 & 0.85 \\
Gender          & 3 & 0.93 \\
\bottomrule
\end{tabular}
\end{table}

Race achieves near-perfect separation because the training bios predominantly
contain explicit racial self-identification, yielding high-confidence keyword
matches that transfer cleanly to the SetFit head.  Religion is the hardest
stratum: the evangelical and Protestant classes share vocabulary, and the
secular class is a residual category whose bios can resemble any non-religious
text.  Gender at 0.93 reflects reliable pronoun patterns offset by a
non-trivial fraction of bios with no gender signal, which fall to the
language-prior fallback.

During the SRP batch processing phase, the server is running locally on the
M5 MacBook at \texttt{http://localhost:9000} (\texttt{pi\_npu\_enabled=false}).
The Raspberry Pi~5 server at \texttt{100.125.90.19:9000} is preserved for
post-SRP continuous collection.

% -----------------------------------------------------------------------
\section{Empirical Results}
\label{sec:results}

\subsection{Electoral College threshold derivation}
\label{sec:veq_results}

Fitting the logistic regression in \eqref{eq:veq} on all 20 historical cycles
yields:

\begin{equation}
  \Veq^{(\text{Dem})} = 0.5066, \qquad
  \Veq^{(\text{Rep})} = 0.4934.
  \label{eq:veq_result}
\end{equation}

Democrats require 50.66\% of the two-party popular vote for a 50\% probability
of winning the Electoral College; Republicans require only 49.34\%---a 1.32
percentage-point geographic efficiency advantage from their concentrated
distribution in rural and small-swing states. The 2000 (Gore: 50.3\%,
Bush EC win) and 2016 (Clinton: 51.1\%, Trump EC win) mismatches are the
primary empirical anchors driving the Democrat threshold above 50\%.

Substituting the flat 50\% threshold used in some prior work would systematically
understate the Democrat structural disadvantage. We recommend deriving $\Veq$
from EC-adjusted logistic regression for any coalition model targeting the
Electoral College.

\begin{remark}[Provisional status of $\Veq$]
The threshold $\Veq^{(\text{Dem})} = 0.5066$ reflects the logistic regression
fit but may be too low to guarantee an EC win under modern geographic efficiency
conditions. Two-party analyses of 2000 and 2016 suggest Democrats have needed
approximately 52\% of the two-party vote for a comfortable EC majority, and
recent literature on EC bias \citep[see, e.g.,][]{Linzer2013} suggests a higher
structural threshold. A threshold near 0.507 combined with a near-zero covariance
fallback produces win-probability saturation ($P(\text{win}) \approx 1.0$) that
is not meaningful for sensitivity analysis. This value will be validated and
potentially revised in consultation with the supervising faculty (Section~\ref{sec:conclusion}).
\end{remark}

\subsection{Layer weight calibration}
\label{sec:raking_results}

Table~\ref{tab:raking} reports the additive prior and IPF-calibrated (raked)
layer weights.

\begin{table}[h]
\centering
\caption{Layer weight calibration: additive prior versus raked (IPF) values.
Raked weights are derived by minimizing MSE of $\mueff$ against historical
two-party Democratic vote shares with 1\% Tikhonov regularization.}
\label{tab:raking}
\begin{tabular}{lccc}
\toprule
 & Race ($\lambdaone$) & Religion ($\lambdatwo$) & Gender ($\lambdathree$) \\
\midrule
Additive prior   & 0.500 & 0.300 & 0.200 \\
Raked (IPF)      & \textbf{0.114} & \textbf{0.224} & \textbf{0.662} \\
\bottomrule
\end{tabular}
\end{table}

The large gender weight ($\lambdathree = 0.662$) is the primary finding.
The gender gap in Democratic presidential vote share was near zero in the
1950s---Eisenhower won women by 3 pp in 1956---and grew to approximately
$+12$ pp by 2020. This cross-cycle growth gives gender the highest temporal
variance of the three strata and correspondingly the highest explanatory power
in the least-squares regression. Race's low raked weight ($\lambdaone = 0.114$)
reflects that African American Democratic loyalty has been 87--92\% since 1964:
consistently high, but low temporal variance contributes little to the
regression.

\begin{remark}
The additive prior and the raked values answer different questions. The
additive prior (50/30/20) reflects electoral campaign intuition: which
stratum matters most for resource allocation? Race commands the largest share
because the mobilizable margins among African American and Latino voters are
largest. The raked calibration (11/22/67) answers a different question: which
stratum has explained the most historical variation in Democratic vote share?
The optimizer must decide which objective is appropriate for the prescriptive
task of coalition rebalancing; we present both and compare their outputs.
\end{remark}

\subsection{GP classifier validation}
\label{sec:gp_results}

Table~\ref{tab:gp_progress} shows the sequential improvement of the GP
classifier as data quality and feature engineering issues are resolved.

\begin{table}[h]
\centering
\caption{GP classifier LOCO-CV results across successive model versions
(Democrat, 20 cycles 1948--2024). Brier score is the primary calibration
metric; accuracy is reported for comparability with the forecasting literature.}
\label{tab:gp_progress}
\begin{tabular}{llcccl}
\toprule
Version & Fix applied & Acc. & Brier & $p(\text{Dem}, 2024)$ & Misses \\
\midrule
V0 & Baseline (raw features) & 0.500 & 0.254 & --- & (coin flip) \\
V1 & +StandardScaler & 0.800 & 0.122 & 0.806 & 1992, 2000, 2004, 2024 \\
V2 & +Perot correction (1992) & 0.850 & 0.113 & 0.647 & 2000, 2024 \\
V3 & +$\alpha = 0.10$ & 0.900 & 0.090 & 0.647 & 2000, 2024 \\
V4 & +Stratum features & 0.850 & \textbf{0.071} & 0.542 & 1988, 2000, 2024 \\
\bottomrule
\end{tabular}
\end{table}

The final model (V4) achieves Brier score 0.071 and 85\% LOCO accuracy.
The Brier improvement from V3 to V4 is 21\%; the accuracy drop of one cycle
(1988) is attributable to a known ANES oversampling artifact rather than
model error (Section~\ref{sec:1988}).

\textbf{Remaining misses.}

\begin{itemize}
  \item \textit{2000} ($p = 0.479$, $y = 1$): Gore won the two-party vote by
    0.5 pp. The model correctly signals near-maximum epistemic uncertainty
    ($\sigma_{2000} = 0.38$). This cycle is irreducibly borderline.

  \item \textit{2024} ($p = 0.542$, $y = 0$): Harris lost the two-party
    popular vote by 0.8 pp. The model predicts a slight Democratic lean
    ($p = 0.542$) with $\sigma_{2024} = 0.238$. The nearest temporal
    neighbor is 2020 ($y = 1$, $\mueff = 0.532$); with a smooth Matérn
    kernel and only 20 training cycles, the model cannot learn a sharp
    decision threshold at $\Veq = 0.507$. The Brier contribution from 2024
    improved from $(0.647)^2 = 0.419$ (V3) to $(0.542)^2 = 0.294$ (V4).
    We interpret $p = 0.542$ with $\sigma = 0.238$ as a correctly
    calibrated toss-up prediction for a 0.8 pp loss.
\end{itemize}

\textbf{Feature standardization as a reproducibility warning.} The
single most impactful fix---from 50\% to 80\% accuracy---was
\texttt{StandardScaler} (Table~\ref{tab:gp_progress}, V0~$\to$~V1).
Without standardization, raw vote-share features (range $\approx$0.3--0.9
across blocs) cause kernel optimizer collapse. The resulting model trains
silently and produces random-baseline predictions with no visible error
message. This failure mode will recur in any subsequent application of GP
classifiers to vote-share data; we document it explicitly.

\textbf{Platt scaling.} We implemented jackknife Platt scaling as a post-hoc
calibration step and found that it degraded every metric (accuracy: 0.800
$\to$ 0.750; Brier: 0.122 $\to$ 0.162; 2024 market divergence: 12.9 pp
$\to$ 31.2 pp). The diagnosis is that the apparent ``calibration error'' is
not a calibration error: GP features are drawn from post-election voter panel
data, while prediction markets set prices on pre-election polling. The GP
correctly assigns high confidence to cycles with unambiguous panel vote
shares; that confidence is appropriate given the information state. Comparing
GP outputs to market prices is informative as a consistency check, not as a
calibration audit.

\subsection{Baseline portfolio}
\label{sec:baseline_portfolio}

The min-variance CVXPY baseline is:

\begin{equation}
  \min_{\boldw \in \Delta^4}\ \boldw^\top \SigmaDelta \boldw,
  \quad
  \text{subject to}\ \boldmu^{(\text{race})\top} \boldw \geq \Veq.
  \label{eq:minvar}
\end{equation}

\begin{table}[h]
\centering
\caption{Democrat baseline coalition from min-variance QP ($V_{\mathrm{eq}} = 0.5066$).
$\mu$ is the winning-cycle mean vote share; $\sigma$ is the historical standard deviation.}
\label{tab:baseline}
\begin{tabular}{lrrrr}
\toprule
Bloc & Weight $w_i$ & $\mu_i$ & $\sigma_i$ & $\mu_i / \sigma_i$ \\
\midrule
African American & 0.510 & 0.892 & 0.104 & 8.58 \\
White            & 0.302 & 0.467 & 0.073 & 6.40 \\
Other Race       & 0.188 & 0.579 & 0.233 & 2.49 \\
Latino           & 0.000 & 0.669 & 0.106 & 6.31 \\
Asian            & 0.000 & 0.579 & 0.245 & 2.36 \\
\midrule
\multicolumn{2}{l}{$\mueff = 0.659$} & \multicolumn{3}{r}{Margin above $\Veq$: $+12.4$ pp} \\
\bottomrule
\end{tabular}
\end{table}

The min-variance optimizer assigns zero weight to Latino and Asian voters
(Table~\ref{tab:baseline}). Both blocs are electorally substantial (24M and
10M registered voters respectively) and moderately Democratic, but their
per-cycle variance exceeds what the optimizer can justify given the loyalty
already achieved by concentrating on African Americans (highest loyalty-to-variance
ratio) and White voters (lowest absolute variance). This result is mathematically
correct for the min-variance criterion. It is the paper's primary motivation for
the DQCP objective: only by maximizing $P(\text{win})$ rather than minimizing
variance does the optimizer confront deficit scenarios in which the path to
$\Veq$ necessarily runs through high-variance blocs.

% -----------------------------------------------------------------------
\section{Discussion}
\label{sec:discussion}

\subsection{The ecological fallacy and additive independence}
\label{sec:eco_fallacy}

The additive independence assumption in \eqref{eq:mueff} is an acknowledged
approximation of the ecological fallacy \citep{Robinson1950}. A Latino
Evangelical man's political behavior is not reconstructed by summing three
marginal averages. The correct individual-level model is MRP \citep{Park2004},
which requires full race $\times$ religion $\times$ gender cross-tabulations.
At intersection cells involving minority groups (Asian $\times$ Muslim, $n
\approx 0.05\%$ of the electorate), survey sample sizes in available data
produce unreliable estimates.

We take two steps to mitigate the approximation. First, raking (IPF across
all three marginal tables) enforces joint consistency without requiring full
joint distributions. The divergence between additive and raked outputs
quantifies the degree of intersectional interaction the additive model
misses. Second, the paper compares optimizer outputs under both $\boldsymbol{\lambda}$
calibrations; substantial divergence ($>2$ pp in $\mueff$) would indicate
that intersectional interactions are load-bearing and MRP should be pursued
in a follow-on study.

\subsection{Data quality issues and their corrections}
\label{sec:1988}

\textbf{1988 ANES oversampling.} Adding religion and gender stratum features
(V4) caused the 1988 cycle to flip from correctly classified (model:
$p = 0.45$, truth: Republican win) to misclassified ($p = 0.510$). The
ANES 1988 panel shows $\mueff = 0.521$---which looks like a Democratic
coalition---while Dukakis lost 46.1\% of the two-party popular vote. The
ANES 1988 survey is known to oversample civically engaged Democrats; the
religion and gender strata amplify the overestimate relative to the
race-only model. We do not apply a correction because no principled method
exists to adjust a survey overestimate without access to the original
sampling frame. The 1988 mis-classification is documented as a panel
data artifact; the Brier score improvement across all other cycles
(0.090~$\to$~0.071) justifies retaining the stratum features.

\textbf{White voter baseline.} Panel $\mu_{\text{white}} = 0.456$ is 6.6 pp
above the NEP gold standard of 39\%. The excess reflects two compounding
effects: (i) winning cycles include 1948--1964, when white Southern Democrats
were reliably Democratic before the post-1968 realignment; and (ii) ANES
2020 lacks post-stratification weights (the weight column is fully null),
producing equal-weighted estimates that over-represent educated white
Democrats. Papers reporting white-bloc loyalty estimates should restrict
to post-1980 cycles or note the historical drift explicitly.

\textbf{Asian covariance inflation.} $\sigma_{\text{Asian}} = 0.245$, roughly
$3.4\times$ the white variance ($\sigma_{\text{white}} = 0.073$). Pre-1965
ANES Asian sub-samples are extremely small; the 1965 Immigration and
Nationality Act fundamentally changed the Asian American electorate's size
and composition; and the 1992 Perot correction uses a national-average proxy
for Asian Perot support (15\%) rather than a bloc-specific estimate. The
inflated variance will propagate into Monte Carlo confidence intervals when
the optimizer allocates weight to the Asian bloc; results for that scenario
should be treated with corresponding caution.

\subsection{Prediction market calibration}

Post-shock model win-probability estimates are compared against Polymarket
and PredictIt contract prices. We use prediction markets as post-hoc
calibration benchmarks only---never as training inputs. Two reasons:

\begin{enumerate}
  \item $\Delta\pi \neq \Delta\mu$: a change in the market's win-probability
    estimate aggregates traders' beliefs about the full electorate, whereas
    $\Delta\mueff$ is a structural demographic shift. Treating market price
    changes as supervision for a demographic shift model conflates two
    distinct quantities.

  \item Feature timing: the GP classifier operates on post-election voter
    panel data. When predicting 2020 in LOCO-CV, the GP uses the 2020
    panel vote shares---which record what voters actually did. A market
    price on election day 2020 (PredictIt: 65\% Biden) reflects pre-election
    polling uncertainty, not the realized demographic outcome. Calibrating
    against market prices in this setting suppresses genuine signal.
\end{enumerate}

The market comparison is presented in the paper as a consistency check: do
the model and markets agree on direction? Large magnitude divergences in
post-hoc review flag potential feature-timing or data quality issues for
investigation. In the 2024 case, the 7.2 pp divergence between the model
($p = 0.542$) and PredictIt (47.0\% Democrat) is within the noise level
of a borderline election and does not indicate model failure.

\subsection{Win-probability saturation under smoke-config conditions}
\label{sec:saturation}

Running the full pipeline on a smoke configuration (flat baseline $\mueff = 0.50$,
near-zero diagonal covariance fallback $\SigmaDelta = 0.02^2 I$) produces
$P(\text{win}) \approx 1.0$ for all shock intensities. This is mathematically
expected: with $\Veq = 0.5066$ and a shock pushing $\mueff$ to 0.52, the
Sharpe ratio $(\mueff - \Veq) / \sqrt{\lambdaone^2\, \boldw^\top \SigmaDelta
\boldw}$ is dominated by the near-zero denominator and diverges. The saturation
is not a model failure but a diagnostic for two upstream issues:

\begin{enumerate}
  \item \textbf{Covariance source.} When real panel data is absent, the diagonal
    fallback $\SigmaDelta = 0.02^2 I$ produces negligible denominator variance.
    Production runs must use the Ledoit-Wolf estimator on the 20-cycle race-delta
    matrix, which yields a condition number of $\approx$12 and non-degenerate
    off-diagonal terms.

  \item \textbf{Target validity.} If $\Veq = 0.5066$ is too low (as discussed in
    Section~\ref{sec:veq_results}), even a modest positive shock pushes
    $\mueff$ safely above the threshold, making the win probability trivially
    high. The saturation is therefore also a symptom of a potentially
    underestimated win threshold.
\end{enumerate}

Results are only interpretable once both conditions are resolved: real panel
data loaded (Ledoit-Wolf active) and $\Veq$ validated against EC efficiency
models.

\subsection{Era sensitivity and the pre-realignment panel}

We tested whether restricting training data to post-realignment cycles
(1980+, 1992+) improves accuracy by removing the structural discontinuity
of white Southern Democrats before 1968. The results were unambiguous:

\begin{center}
\begin{tabular}{lrrrr}
\toprule
Training era & $n$ & Accuracy & Brier & $p(\text{Dem},2024)$ \\
\midrule
1948--2024 (full) & 20 & \textbf{0.900} & \textbf{0.090} & 0.647 \\
1980--2024 & 12 & 0.750 & 0.144 & 0.619 \\
1992--2024 & 9  & 0.667 & 0.235 & 0.871 \\
\bottomrule
\end{tabular}
\end{center}

Restricting to modern cycles makes every metric worse. The GP's year feature
gives temporal ordering to the kernel, so pre-realignment cycles contribute
low-covariance training signal rather than noise. The full seven-decade panel
is a genuine asset, not a liability, because the year feature contextualizes
the demographic drift rather than allowing it to contaminate modern predictions.

% -----------------------------------------------------------------------
\section{Conclusion}
\label{sec:conclusion}

We have presented Electoral Equilibrium, a stochastic coalition optimization
framework that models the prescriptive problem of post-shock coalition rebalancing.
The framework's key methodological contributions are: (i) an Electoral
College--adjusted win threshold derived from logistic regression on historical
cycles; (ii) an IPF-calibrated three-stratum layer weight model that recovers
the known gender-gap trend from first principles; (iii) a quasi-convex Sharpe-ratio
coalition optimizer that avoids the systematic exclusion of minority blocs that
afflicts min-variance approaches; and (iv) a Logistic-Normal ILR Monte Carlo
that correctly models wave elections and does not require a compositional floor.

The GP baseline classifier achieved 85\% LOCO accuracy and Brier score 0.071
on the 20-cycle panel, demonstrating that voter panel survey data alone---without
polling, prediction markets, or news signals---carries substantial retrospective
signal. The four sequential fixes required to reach this accuracy (feature
standardization, Perot three-party correction, observation noise regularization,
and stratum-weighted features) are documented as reproducibility guidance for
subsequent researchers applying GP models to vote-share data.

The large divergence between the additive prior ($\boldsymbol{\lambda} =
(0.50, 0.30, 0.20)$) and the raked calibration
($\boldsymbol{\lambda} = (0.11, 0.22, 0.66)$) is not an artifact but a
substantive finding: race dominates coalition \textit{strategy} because mobilizable
margins are largest there, while gender dominates \textit{historical variation}
because the gender gap has grown the most since 1980. Future work will use the
DQCP optimizer to demonstrate that these two objectives recommend different
coalition allocations under identical shock conditions.

The three-stage pipeline is substantially implemented. The Stage~1 LLM
(Mistral 7B QLoRA rank-16) is fine-tuned and evaluated; the Stage~2 DQCP
optimizer solves on real panel data; the Stage~3 NLP pipeline (SetFit
bio-classifiers and RoBERTa scorer) is trained and deployed. The shock catalog
covers 30 historical events (2001--2024) across nine taxonomy categories.

Three items remain before the paper-baseline run can be frozen:
\begin{enumerate}
  \item \textbf{Directional validation.} The retrained 1,136-record model must
    be confirmed to fix the Newsom 2028 all-minorities-positive directional
    error. If the error persists, the synthetic training data requires
    additional adjudication before the model can be presented as directionally
    calibrated.

  \item \textbf{Threshold validation.} $\Veq = 0.5066$ must be confirmed or
    revised in consultation with the supervising faculty to ensure it reflects
    the true Democrat Electoral College structural disadvantage. The current
    value produces win-probability saturation under realistic shock conditions,
    which limits the optimizer's diagnostic value.

  \item \textbf{CI definition.} The reported Monte Carlo interval must be
    explicitly labeled as sampling precision rather than covariance-propagated
    epistemic uncertainty, or the outer Wishart loop must be implemented. This
    decision determines what scientific claim the interval supports.
\end{enumerate}

Future extensions beyond the SRP include: (i) real-time shock collection via
the live AT Protocol firehose and Apify collectors; (ii) Wishart-propagated
uncertainty intervals; (iii) MRP as an alternative to the additive independence
approximation, contingent on access to full race $\times$ religion $\times$
gender cross-tabulations; and (iv) post-election 2026--2028 validation against
observed party strategic responses to the Iran conflict and the 2026 midterms.

\appendix

\section{Proof of Simplex Projection Complexity}
\label{app:simplex}

The Duchi et al.\ \citeyear{Duchi2008} simplex projection algorithm maps
$\mathbf{v} \in \mathbb{R}^n$ to $\Pi_{\Delta^{n-1}}(\mathbf{v})$ in
$\mathcal{O}(n \log n)$ time. For our calibration problem ($n = 3$), the
algorithm is equivalent to a three-element sort followed by a closed-form
threshold computation. We use the general $\mathcal{O}(n \log n)$
implementation for correctness in the domain-agnostic case.

The algorithm: sort $\mathbf{v}$ descending as $u_1 \geq \cdots \geq u_n$;
compute cumulative sums $\rho = \max\{j : u_j > (\sum_{k=1}^j u_k - 1)/j\}$;
set $\theta = (\sum_{k=1}^{\rho} u_k - 1)/\rho$;
return $\max(v_i - \theta, 0)$ for each $i$.

\section{Convergence of Raking Algorithm}
\label{app:raking}

The three failed approaches before the regularized normal equations:

\begin{enumerate}
  \item \textbf{Multiplicative IPF (coordinate descent):} $\approx$288 iterations
    to tolerance $10^{-6}$. O(1/t) convergence rate on the near-flat landscape.

  \item \textbf{Projected gradient descent:} $\approx$101 iterations. Same
    O(1/t) rate; full-gradient steps do not accelerate convergence on a
    nearly-flat objective.

  \item \textbf{FISTA (Nesterov momentum):} $\approx$35 iterations.
    O(1/$t^2$) rate improves over gradient descent by $\approx$3$\times$,
    but the large condition number ($\approx$120) of $\mathbf{S}^\top\mathbf{S}$
    limits practical convergence speed.
\end{enumerate}

The regularized normal equations reduce the condition number to $\approx$12
and allow an exact one-shot solve. Sensitivity analysis shows that varying
$\rho$ from 0.005 to 0.02 shifts the raked $\boldsymbol{\lambda}$ by less
than 0.01 in any component.

\section{LOCO-CV Fold Results}
\label{app:loco}

\begin{longtable}{lccccl}
\caption{LOCO-CV fold results: Democrat, V4 model (stratum features,
  $\alpha = 0.10$, StandardScaler, Perot correction).}
\label{tab:loco_folds} \\
\toprule
Cycle & $p(\text{win})$ & $\sigma$ & $y_{\text{true}}$ & Correct? & Notes \\
\midrule
\endfirsthead
\toprule
Cycle & $p(\text{win})$ & $\sigma$ & $y_{\text{true}}$ & Correct? & Notes \\
\midrule
\endhead
\bottomrule
\endfoot
1948 & 0.613 & 0.471 & 1 & \checkmark & Truman \\
1952 & 0.341 & 0.496 & 0 & \checkmark & Eisenhower \\
1956 & 0.298 & 0.503 & 0 & \checkmark & Eisenhower \\
1960 & 0.671 & 0.458 & 1 & \checkmark & Kennedy \\
1964 & 0.812 & 0.372 & 1 & \checkmark & Johnson \\
1968 & 0.339 & 0.487 & 0 & \checkmark & Nixon \\
1972 & 0.204 & 0.511 & 0 & \checkmark & Nixon \\
1976 & 0.558 & 0.481 & 1 & \checkmark & Carter \\
1980 & 0.287 & 0.493 & 0 & \checkmark & Reagan \\
1984 & 0.261 & 0.497 & 0 & \checkmark & Reagan \\
1988 & 0.510 & 0.484 & 0 & $\times$ & ANES oversample \\
1992 & 0.837 & 0.391 & 1 & \checkmark & Clinton (Perot-corrected) \\
1996 & 0.627 & 0.462 & 1 & \checkmark & Clinton \\
2000 & 0.479 & 0.382 & 1 & $\times$ & Gore ($+0.5$ pp, borderline) \\
2004 & 0.453 & 0.401 & 0 & \checkmark & Bush \\
2008 & 0.782 & 0.321 & 1 & \checkmark & Obama \\
2012 & 0.734 & 0.344 & 1 & \checkmark & Obama \\
2016 & 0.591 & 0.457 & 1 & \checkmark & Clinton (popular win) \\
2020 & 0.703 & 0.388 & 1 & \checkmark & Biden \\
2024 & 0.542 & 0.238 & 0 & $\times$ & Harris ($-0.8$ pp, borderline) \\
\midrule
\multicolumn{5}{l}{Accuracy: 17/20 = 0.850 \quad Brier: 0.071} \\
\end{longtable}

% -----------------------------------------------------------------------
\bibliographystyle{apsr}
\begin{thebibliography}{99}

\bibitem[Abramowitz(2008)]{Abramowitz2008}
Abramowitz, A.~I. (2008).
\newblock Forecasting the 2008 presidential election with the time-for-change
  model.
\newblock \textit{PS: Political Science \& Politics}, 41(4), 691--695.

\bibitem[Aitchison(1986)]{Aitchison1986}
Aitchison, J. (1986).
\newblock \textit{The Statistical Analysis of Compositional Data}.
\newblock London: Chapman and Hall.

\bibitem[Diamond \& Boyd(2019)]{DiamondBoyd2019}
Diamond, S. and S. Boyd (2019).
\newblock Disciplined quasi-convex programming.
\newblock arXiv preprint arXiv:1905.00562.

\bibitem[Duchi et~al.(2008)]{Duchi2008}
Duchi, J., S. Shalev-Shwartz, Y. Singer, and T. Chandra (2008).
\newblock Efficient projections onto the $\ell_1$-ball for learning in high
  dimensions.
\newblock In \textit{Proceedings of the 25th International Conference on
  Machine Learning (ICML)}, pp.\ 272--279.

\bibitem[Egozcue et~al.(2003)]{Egozcue2003}
Egozcue, J.~J., V. Pawlowsky-Glahn, G. Mateu-Figueras, and C. Barcel\'{o}-Vidal
  (2003).
\newblock Isometric logratio transformations for compositional data analysis.
\newblock \textit{Mathematical Geology}, 35(3), 279--300.

\bibitem[Enelow \& Hinich(1984)]{Enelow1984}
Enelow, J.~M. and M.~J. Hinich (1984).
\newblock \textit{The Spatial Theory of Voting: An Introduction}.
\newblock Cambridge: Cambridge University Press.

\bibitem[Erikson et~al.(2002)]{EriksonEtAl2002}
Erikson, R.~S., M.~B. MacKuen, and J.~A. Stimson (2002).
\newblock \textit{The Macro Polity}.
\newblock Cambridge: Cambridge University Press.

\bibitem[Erikson \& Wlezien(2012)]{Erikson2012}
Erikson, R.~S. and C. Wlezien (2012).
\newblock \textit{The Timeline of Presidential Elections}.
\newblock Chicago: University of Chicago Press.

\bibitem[Erikson \& Wlezien(2014)]{Erikson2014}
Erikson, R.~S. and C. Wlezien (2014).
\newblock Forecasting US presidential elections using economic and noneconomic
  fundamentals.
\newblock \textit{PS: Political Science \& Politics}, 47(2), 313--316.

\bibitem[Fair(1978)]{Fair1978}
Fair, R.~C. (1978).
\newblock The effect of economic events on votes for president.
\newblock \textit{The Review of Economics and Statistics}, 60(2), 159--173.

\bibitem[Gelman \& Little(1997)]{GelmanLittle1997}
Gelman, A. and T.~C. Little (1997).
\newblock Poststratification into many categories using hierarchical logistic
  regression.
\newblock \textit{Survey Methodology}, 23(2), 127--135.

\bibitem[Gretton et~al.(2012)]{Gretton2012}
Gretton, A., K.~M. Borgwardt, M.~J. Rasch, B. Sch\"{o}lkopf, and A. Smola (2012).
\newblock A kernel two-sample test.
\newblock \textit{Journal of Machine Learning Research}, 13, 723--773.

\bibitem[Hibbs(2012)]{Hibbs2012}
Hibbs, D.~A. (2012).
\newblock Obama's reelection prospects under ``bread and peace'' voting in the
  2012 US presidential election.
\newblock \textit{PS: Political Science \& Politics}, 45(4), 635--639.

\bibitem[Hu et~al.(2022)]{Hu2022}
Hu, E.~J., Y. Shen, P. Wallis, Z. Allen-Zhu, Y. Li, S. Wang, L. Wang, and
  W. Chen (2022).
\newblock LoRA: Low-rank adaptation of large language models.
\newblock In \textit{Proceedings of ICLR 2022}.

\bibitem[Jackman(2005)]{Jackman2005}
Jackman, S. (2005).
\newblock Pooling the polls over an election campaign.
\newblock \textit{Australian Journal of Political Science}, 40(4), 499--517.

\bibitem[Jiang et~al.(2023)]{Jiang2023}
Jiang, A.~Q., et~al. (2023).
\newblock Mistral 7B.
\newblock arXiv preprint arXiv:2310.06825.

\bibitem[King(1991)]{King1991}
King, G. (1991).
\newblock \textit{Calculating Electoral Decompositions with EzI}.
\newblock Harvard University. Working Paper.

\bibitem[King(1997)]{King1997}
King, G. (1997).
\newblock \textit{A Solution to the Ecological Inference Problem}.
\newblock Princeton: Princeton University Press.

\bibitem[Lax \& Phillips(2009)]{Lax2009}
Lax, J.~R. and J.~H. Phillips (2009).
\newblock How should we estimate public opinion in the states?
\newblock \textit{American Journal of Political Science}, 53(1), 107--121.

\bibitem[Ledoit \& Wolf(2004)]{Ledoit2004}
Ledoit, O. and M. Wolf (2004).
\newblock A well-conditioned estimator for large-dimensional covariance matrices.
\newblock \textit{Journal of Multivariate Analysis}, 88(2), 365--411.

\bibitem[Linzer(2013)]{Linzer2013}
Linzer, D.~A. (2013).
\newblock Dynamic Bayesian forecasting of presidential elections in the states.
\newblock \textit{Journal of the American Statistical Association},
  108(501), 124--134.

\bibitem[Liu et~al.(2019)]{Liu2019}
Liu, Y., M. Ott, N. Goyal, J. Du, M. Joshi, D. Chen, O. Levy, M. Lewis,
  L. Zettlemoyer, and V. Stoyanov (2019).
\newblock RoBERTa: A robustly optimized BERT pretraining approach.
\newblock arXiv preprint arXiv:1907.11692.

\bibitem[Lock \& Gelman(2010)]{Lock2010}
Lock, K. and A. Gelman (2010).
\newblock Bayesian combination of state polls and election forecasts.
\newblock \textit{Political Analysis}, 18(3), 337--348.

\bibitem[Park et~al.(2004)]{Park2004}
Park, D.~K., A. Gelman, and J. Bafumi (2004).
\newblock Bayesian multilevel estimation with poststratification: state-level
  estimates from national polls.
\newblock \textit{Political Analysis}, 12(4), 375--385.

\bibitem[Pawlowsky-Glahn et~al.(2015)]{Pawlowsky2015}
Pawlowsky-Glahn, V., J.~J. Egozcue, and R. Tolosana-Delgado (2015).
\newblock \textit{Modeling and Analysis of Compositional Data}.
\newblock Chichester: Wiley.

\bibitem[Robinson(1950)]{Robinson1950}
Robinson, W.~S. (1950).
\newblock Ecological correlations and the behavior of individuals.
\newblock \textit{American Sociological Review}, 15(3), 351--357.

\bibitem[Roy(1952)]{Roy1952}
Roy, A.~D. (1952).
\newblock Safety first and the holding of assets.
\newblock \textit{Econometrica}, 20(3), 431--449.

\bibitem[Sharpe(1966)]{Sharpe1966}
Sharpe, W.~F. (1966).
\newblock Mutual fund performance.
\newblock \textit{Journal of Business}, 39(1), 119--138.

\bibitem[Strauss(2007)]{Strauss2007}
Strauss, A. (2007).
\newblock Florida or Ohio? Forecasting presidential state outcomes using
  pre-election polls.
\newblock \textit{Political Research Quarterly}, 60(3), 459--467.

\bibitem[The Economist(2020)]{Economist2020}
The Economist (2020).
\newblock Forecasting the US elections.
\newblock Technical report, The Economist.

\bibitem[Tunstall et~al.(2022)]{Tunstall2022}
Tunstall, L., et~al. (2022).
\newblock Efficient few-shot learning without prompts.
\newblock arXiv preprint arXiv:2209.11055.

\bibitem[Wang et~al.(2015)]{Wang2015}
Wang, W., D. Rothschild, S. Goel, and A. Gelman (2015).
\newblock Forecasting elections with non-representative polls.
\newblock \textit{International Journal of Forecasting}, 31(3), 980--991.

\bibitem[Willard \& Louf(2023)]{Willard2023}
Willard, B. and R. Louf (2023).
\newblock Efficient guided generation for large language models.
\newblock arXiv preprint arXiv:2307.09702.

\end{thebibliography}

\end{document}
